\documentclass[10pt,a4paper]{amsart}
\usepackage[margin=19mm]{geometry}
\usepackage{amsmath,amssymb,mathtools}
\usepackage[scr=rsfs,cal=euler]{mathalfa}
\usepackage{xfrac}
\usepackage[unicode,hidelinks,bookmarks=false]{hyperref}
\newcommand{\ie}{i.e.\ }
\newcommand{\cf}{cf.\ }
\newcommand{\resp}{respectively\ }
\newcommand{\Subsection}{\subsection}
\providecommand{\nobreakdash}{\nobreak\hbox{-}}
\setlength{\emergencystretch}{2em}
\multlinegap0pt
\usepackage{bm}
\usepackage{bbm}
\usepackage{dsfont}
\usepackage{xspace}
\usepackage{cancel}
\usepackage{mathtools}
\usepackage{amsthm}
\makeatletter
\@ifundefined{proposition}{\newtheorem{proposition}{Proposition}}{}
\@ifundefined{theorem}{\newtheorem{theorem}{Theorem}}{}
\makeatother
\DeclarePairedDelimiter{\abs}{\lvert}{\rvert}
\newcommand{\bbC}{\mathbb{C}}
\newcommand{\bbN}{\mathbb{N}}
\newcommand{\calO}{\mathcal{O}}
\newcommand{\dd}{\,\mathrm{d}}
\newcommand{\rme}{\mathrm{e}}
\newcommand{\rmi}{\mathrm{i}}
\renewcommand{\Re}{\operatorname{Re}}
\title[On a dense set of functions determined by sampled Gabor magn]{On a dense set of functions determined by sampled Gabor magnitude}
\author{Matthias Wellershoff}
\date{Source: arXiv:2507.14556v2 (CC BY 4.0)}
\begin{document}
\maketitle

\noindent\emph{Source and licence.} On a dense set of functions determined by sampled Gabor magnitude. Version arXiv:2507.14556v2 (CC BY 4.0), distributed under the Creative Commons Attribution 4.0 International licence, as recorded in the arXiv licence metadata for this exact version and shown on the abstract page https://arxiv.org/abs/2507.14556v2. Full rights notice: licenses/arxiv-2507.14556-CC-BY-4.0.txt.\par\smallskip

\noindent\textit{Editorial scope.} Counted unit: the Proposition whose source label is \texttt{prop:PolynomialProposition} in the article (source0.tex lines 573--579, source label \texttt{prop:PolynomialProposition}), delivered together with the article's own complete original proof of it (source0.tex lines 581--639, one \texttt{proof} environment of 59 lines). No mathematics is written, repaired or re-derived: statement and proof are the article's own text line for line. Required context retained uncounted: eq:focklikecondition (lines 314--316); thm:Earl (lines 325--334). Every definition, display and label of those lines is retained unchanged. Omitted from this package and not redistributed: 1--313, 317--324, 335--572, 580--580, 640--766. Nothing inside a retained excerpt is omitted or altered: every retained body is delivered exactly as the article's lines are written, so no omission receipt of any kind accompanies this package. Citations: the retained excerpts carry no citation command at all, so no bibliography entry of the article is transcribed and no reference is redistributed. Reference adaptation: none was needed and none was made; every reference inside the delivered unit resolves inside it, so no unresolved reference or citation is produced. Printed slips kept literal: no printed word, symbol, hypothesis or number of the counted statement or of its proof was altered. Layout and class: the article's own class is \texttt{article} with packages including amsfonts, amsthm, amssymb, bbm, hyperref, mathtools, multicol, tikz, pgfplots, xcolor, xurl; the wrapper is the lane's pinned \texttt{amsart} editorial wrapper at 10pt on A4, which restates the theorem-like environments the retained text needs and adds the article's own macro definitions that the retained text needs (\texttt{\textbackslash Re}, \texttt{\textbackslash abs}, \texttt{\textbackslash bbC}, \texttt{\textbackslash bbN}, \texttt{\textbackslash calO}, \texttt{\textbackslash dd}, \texttt{\textbackslash rme}, \texttt{\textbackslash rmi}), with extra wrapper packages (bm, bbm, dsfont, xspace, cancel, mathtools). The delivered numbering is the unit's own. Assets: no retained span contains a figure environment, a graphics-inclusion command, a PDF-page inclusion, a table or a TikZ picture; no image file is copied into this package, the package declares no asset dependency and no image file is redistributed with it. Licence: version arXiv:2507.14556v2 of this article is distributed under the Creative Commons Attribution 4.0 International licence, as recorded in the arXiv licence metadata for this exact version and shown on the abstract page https://arxiv.org/abs/2507.14556v2. A full rights notice is delivered at licenses/arxiv-2507.14556-CC-BY-4.0.txt. Characters: every retained excerpt is pure ASCII, so the delivered engine, XeLaTeX with its default Latin Modern text font, reports no missing character.
    \begin{equation}
        \liminf_{r \to \infty} \frac{\log \max_{z \in C_r} \abs{f(z)}}{r^2} < \infty. \label{eq:focklikecondition}
    \end{equation}

    \begin{theorem}\label{thm:Earl}
        Let $f \in \calO(\bbC)$ be an entire function satisfying condition~\eqref{eq:focklikecondition}. Let $\Lambda = \lbrace \lambda_j \rbrace_{j = 1}^\infty \subset \bbC$ be a uniformly discrete set with lower Beurling density $D^{-}(\Lambda) > 2 \underline \tau / \pi$ such that 
        \begin{equation*}
            \kappa := \limsup_{j \to \infty} \frac{\log \abs{f(\lambda_j)}}{H(\abs{\lambda_j})} < \infty,
        \end{equation*}
        where $H : (0,\infty) \to (0,\infty)$ is non-decreasing and unbounded at infinity such that $r^{-2} H(r)$ is non-increasing and tends to zero as $r \to \infty$. Then, 
        \begin{equation*}
            \limsup_{r \to \infty} \frac{\log \max_{z \in C_r} \abs{f(z)}}{H(r)} \leq \kappa.
        \end{equation*}
    \end{theorem}

    \begin{proposition}\label{prop:PolynomialProposition}
        Let $f \in \calO(\bbC)$ be an entire function satisfying condition~\eqref{eq:focklikecondition} and let $g \in \bbC[z]$ be a polynomial of degree $q \in \bbN_0$. Let $\Lambda \subset \bbC$ be a uniformly discrete set with lower Beurling density $D^{-}(\Lambda) > 2 \underline \tau / \pi$. 
        \begin{enumerate}
            \item If $\abs{f} \lesssim \abs{g}$ on $\Lambda$, then $f \in \bbC[z]$ is a polynomial of degree $q$.
            \item Moreover, $f \sim g$ if and only if $\abs{f} = \abs{g}$ on $\Lambda$.
        \end{enumerate}
    \end{proposition}

    \begin{proof}
        \begin{enumerate}
            \item Since $\abs{f} \lesssim \abs{g}$ on $\Lambda = \lbrace \lambda_j \rbrace_{j = 1}^\infty \subset \bbC$ and $g$ is a polynomial of degree $q \in \bbN_0$, we have 
            \begin{equation*}
                \limsup_{j \to \infty} \frac{\log \abs{f(\lambda_j)}}{\log(1+ \abs{\lambda_j})} \leq \limsup_{j \to \infty} \frac{\log \abs{g(\lambda_j)}}{\log(1+ \abs{\lambda_j})} = q.
            \end{equation*}
            Therefore, Theorem~\ref{thm:Earl} implies that 
            \begin{equation*}
                \limsup_{r \to \infty} \frac{\log \max_{z \in C_r} \abs{f(z)}}{\log(1+r)} \leq q.
            \end{equation*}
            In particular,
            \begin{equation*}
                \abs{f(z)} \lesssim (1 + \abs{z})^{q+1/2}, \qquad z \in \bbC.
            \end{equation*}
            We can now conclude by a standard argument\footnote{The author learned this argument through the answer of Adamski to \url{https://math.stackexchange.com/questions/143468/entire-function-bounded-by-a-polynomial-is-a-polynomial}}: $f$ is an entire function and is therefore given by the convergent power series 
            \begin{equation*}
                f(z) = \sum_{j = 0}^\infty \frac{f^{(j)}(0)}{j!} z^j.
            \end{equation*}
            Additionally, our estimate above together with Cauchy's estimate says that
            \begin{equation*}
                \abs{f^{(j)}(0)} \lesssim \frac{j!}{r^j} \cdot (1 + r)^{q+1/2}, \qquad j \in \bbN_0,~r > 0.
            \end{equation*}
            Letting $r$ go to infinity, we see that $f^{(j)}(0) = 0$ for $j > q$ and thus $f \in \bbC[z]$ is a polynomial of degree $q$ as claimed.

            \item It is clear that $f \sim g$ implies that $\abs{f} = \abs{g}$. So, assume that $\abs{f} = \abs{g}$ on $\Lambda$ and note that item~1 implies that $f$ is a polynomial of degree $q$. By a direct computation,
            \begin{equation*}
                \frac{f(z)}{g(z)} = \frac{a}{b} + \frac{h(z)}{b g(z)}
            \end{equation*}
            for $z \in \mathbb C$ not a zero of $g$, where $a,b \in \mathbb C \setminus \lbrace 0 \rbrace$ are the leading coefficients of $f$ and $g$, respectively, and $h \in \mathbb C [z]$ has degree $< q$. Taking absolute values and using the existence of a sequence $(\lambda_j)_{j = 1}^\infty \subset \Lambda$ tending to infinity, we obtain $\smash{a/b = \rme^{\rmi \alpha}}$ for some $\alpha \in \mathbb R$.

            Suppose by contradiction that $f \not\sim g$ and expand $f/g - \smash{\mathrm{e}^{\mathrm i \alpha}} = h / (bg)$ into a Laurent series
            \begin{equation*}
                \frac{f(z)}{g(z)} - \mathrm e^{\mathrm i \alpha} = \sum_{k = -\infty}^\infty c_k z^{-k}
            \end{equation*}
            outside of a sufficiently large ball. Let $\gamma_r : [0,2\pi] \to \mathbb C$ denote the contour $\gamma_r(t) = r \mathrm e^{\mathrm i t}$ around $B_r$. Since the right-hand side of
            \begin{equation*}
                c_k = \frac{1}{2\pi \mathrm i b} \oint_{\gamma_r} \frac{z^{k-1} h(z)}{g(z)} \dd z = \frac{1}{2\pi b} \int_0^{2\pi} \frac{(r \mathrm e^{\mathrm i t})^k h(r \mathrm e^{\mathrm i t})}{g(r \mathrm e^{\mathrm i t})} \dd t
            \end{equation*}
            tends to zero as $r \to \infty$ when $k \leq 0$, we conclude $c_k = 0$ for all $k \leq 0$, and hence
            \begin{equation*}
                \frac{f(z)}{g(z)} = \mathrm e^{\mathrm i \alpha} + \sum_{k = k_0}^\infty c_k z^{-k},
            \end{equation*}
            for some $k_0 > 0$ with $c_{k_0} \neq 0$ (where we used that $f \not\sim g$ to guarantee that some such non-zero coefficient exists).

            It is immediate that
            \begin{align*}
                \frac{\lvert f(z) \rvert^2}{\lvert g(z) \rvert^2} &= 1 + 2 \Re \left[\mathrm e^{-\mathrm i \alpha} \sum_{k = k_0}^\infty c_k z^{-k}\right] + \left\lvert \sum_{k = k_0}^\infty c_k z^{-k} \right\rvert^2 \\
                &\geq 1 + 2 \Re \left[\mathrm e^{-\mathrm i \alpha} c_{k_0} z^{-k_0}\right] - 2 \sum_{k = k_0+1}^\infty \lvert c_k \rvert \lvert z \rvert^{-k}.
            \end{align*}
            Since $D^-(\Lambda) > 0$, we may find arbitrarily large $\lambda \in \Lambda$ within angle $\theta$ of $\smash{w:=(\mathrm e^{\mathrm i \alpha} \overline c_{k_0})^{-1/k_0}}$, for any fixed $0 \leq \theta < \pi / (2k_0)$. For such $\lambda$,
            \begin{align*}
                1 = \frac{\lvert f(\lambda) \rvert^2}{\lvert g(\lambda) \rvert^2} &\geq 1 + 2 \Re \left[(z \overline w)^{-k_0}\right] - 2 \sum_{k = k_0+1}^\infty \lvert c_k \rvert \lvert \lambda \rvert^{-k} \\
                &\geq 1 + 2 \lvert \lambda \rvert^{-k_0} \lvert c_{k_0} \rvert \cos (k_0 \theta) - 2 \sum_{k = k_0+1}^\infty \lvert c_k \rvert \lvert \lambda \rvert^{-k}.
            \end{align*}
            Once $\lvert \lambda \rvert$ is sufficiently large, the second term dominates the third, yielding $1 > 1$, a contradiction. We conclude $f \sim g$.
        \end{enumerate}

        \vspace{-12pt}
    \end{proof}

\end{document}
